\documentclass[11pt, aspectratio=169]{beamer}
\usepackage[utf8]{inputenc}

\usetheme[color={0, 80, 150}, header]{CambridgeUK}

\begin{document}

    \author{Your Name}
    \title[CambridgeUK]{CambridgeUK Beamer Theme}
    \subtitle{A Complete Command Reference}
    \institute{Institute of Astronomy (IOA)}
    \date{\today}

    % ==========================================
    % TITLE SLIDES
    % ==========================================

    % \titleframe[logo] -- plain title slide
    % logo is optional: path to an institutional logo image
    \titleframe[black_ioa.png]

    % ==========================================
    % THEME SETUP
    % ==========================================

    \section{Theme Setup}
    \sectionframe{Theme Setup}{\texttt{\textbackslash usetheme} options}

    \textframe{Loading the Theme}{
        \texttt{\textbackslash usetheme[options]\{CambridgeUK\}}
        \vspace{0.8em}
        \begin{itemize}
            \item \texttt{color=\{R, G, B\}} --- accent colour in RGB (default: \texttt{0, 80, 150})
            \item \texttt{header} --- show the section name header bar (hidden by default)
            \item \texttt{nav} --- show Beamer navigation symbols (hidden by default)
            \item \texttt{autosec} --- auto-insert a \texttt{\textbackslash sectionframe} at each \texttt{\textbackslash section\{\}}
            \item \texttt{bgimage=\{path\}} --- global background image for all slides
            \item \texttt{bgopacity=\{val\}} --- global background image opacity (default: \texttt{0.3})
            \item \texttt{bgcrop=\{l|p\}} --- global background image crop mode (default: \texttt{l})
        \end{itemize}
    }

    % ==========================================
    % TITLE SLIDES
    % ==========================================

    \section{Title Slides}
    \sectionframe{Title Slides}{\texttt{\textbackslash titleframe} and \texttt{\textbackslash imagetitleframe}}

    \textframe{\texttt{\textbackslash titleframe[logo]}}{
        Produces the plain title slide shown at the start of this document.
        \vspace{0.8em}
        \begin{itemize}
            \item Reads \texttt{\textbackslash title}, \texttt{\textbackslash subtitle}, \texttt{\textbackslash author},
                  \texttt{\textbackslash institute}, and \texttt{\textbackslash date} from the preamble
            \item \texttt{[logo]} is optional --- path to an institutional logo image
            \item Always rendered as a \texttt{[plain]} frame (no header or footer)
        \end{itemize}
    }

    % \imagetitleframe[pos]{crop}{image}
    % pos  -- banner position: north or south (default: south)
    % crop -- l fits by height (landscape); p fits by width (portrait)
    % image -- path to background image
    \imagetitleframe[south]{p}{Earth_From_the_Perspective_of_Artemis_II.jpeg}

    \textframe{\texttt{\textbackslash imagetitleframe[pos]\{crop\}\{image\}}}{
        The previous slide was produced by \texttt{\textbackslash imagetitleframe}.
        \vspace{0.8em}
        \begin{itemize}
            \item \texttt{[pos]} --- translucent banner position: \texttt{north} or \texttt{south} (default: \texttt{south})
            \item \texttt{crop} --- \texttt{l} fits image by height (landscape); \texttt{p} fits by width (portrait)
            \item \texttt{image} --- path to background image
            \item Title, subtitle, and author are drawn from preamble metadata
            \item Always a \texttt{[plain]} frame
        \end{itemize}
    }

    % ==========================================
    % CONTENT SLIDES
    % ==========================================

    \section{Content Slides}
    \sectionframe{Content Slides}{%
        \texttt{\textbackslash textframe},
        \texttt{\textbackslash imageframe},
        \texttt{\textbackslash splitframe}}

    \textframe{\texttt{\textbackslash textframe\{title\}\{content\}}}{
        The standard workhorse slide --- a titled frame with free-form content.
        \vspace{0.8em}
        \begin{itemize}
            \item \texttt{title} --- the frame title
            \item \texttt{content} --- any Beamer content: text, \texttt{itemize}, columns, blocks, etc.
            \item This very slide is a \texttt{\textbackslash textframe}
        \end{itemize}
    }

    % \imageframe[title]{image}[credit]
    % title  -- optional frame title
    % image  -- path to image file; scaled to fit the content area, aspect ratio preserved
    % credit -- optional credit label rendered at the bottom-right corner
    \imageframe[Centred Image Slide: \texttt{\textbackslash imageframe}]{apj460815f8_hr.jpg}[Kopparapu et al. 2013, ApJ, 765, 131]

    \textframe{\texttt{\textbackslash imageframe[title]\{image\}[credit]}}{
        The previous slide was produced by \texttt{\textbackslash imageframe}.
        \vspace{0.8em}
        \begin{itemize}
            \item \texttt{[title]} --- optional frame title; omit for a clean full-content image
            \item \texttt{image} --- path to image file; scales to fit, aspect ratio preserved
            \item \texttt{[credit]} --- optional credit label, rendered via \texttt{\textbackslash credit}
        \end{itemize}
    }

    % \splitframe[width]{align}{title}{image}{content}
    % width -- image column width (default: 0.45\textwidth)
    % align -- l puts image left; r puts image right
    % Use \credit{text} inside the content argument to add a plot citation
    \splitframe[0.5\textwidth]{l}{Split Frame: Image Left --- \texttt{\textbackslash splitframe}}{apj460815f8_hr.jpg}{
        Places an image and text side by side.
        \vspace{0.5em}
        \begin{itemize}
            \item \texttt{[width]} --- image column width (default: \texttt{0.45\textbackslash textwidth})
            \item \texttt{align} --- \texttt{l} image left, \texttt{r} image right
            \item \texttt{title}, \texttt{image}, \texttt{content} --- as labelled
            \item Use \texttt{\textbackslash credit\{text\}} in the content for plot citations
        \end{itemize}
        \credit{Kopparapu et al. 2013, ApJ, 765, 131}
    }

    \splitframe[0.5\textwidth]{r}{Split Frame: Image Right}{apj460815f8_hr.jpg}{
        Same command with \texttt{align = r}.
        \vspace{0.5em}
        \begin{itemize}
            \item Image moves to the right column
            \item Text occupies the remaining width on the left
        \end{itemize}
    }

    % ==========================================
    % CINEMATIC LAYOUTS
    % ==========================================

    \section{Cinematic Layouts}
    \sectionframe{Cinematic Layouts}{\texttt{\textbackslash coverframe} and \texttt{\textbackslash coversplitframe}}

    % \coverframe[crop]{image}[credit]
    % crop   -- l fits by height (landscape); p fits by width (portrait); default: l
    % image  -- path to background image
    % credit -- optional; shown as a small dark pill at the bottom-right corner
    \coverframe[p]{Earth_From_the_Perspective_of_Artemis_II.jpeg}[Credit: NASA/Artemis II]

    \textframe{\texttt{\textbackslash coverframe[crop]\{image\}[credit]}}{
        The previous slide was a \texttt{\textbackslash coverframe} with a credit label.
        \vspace{0.8em}
        \begin{itemize}
            \item \texttt{[crop]} --- \texttt{l} fits by height; \texttt{p} fits by width (default: \texttt{l})
            \item \texttt{[credit]} --- optional image credit, rendered by \texttt{\textbackslash credit}
            \item Always a \texttt{[plain]} frame --- no header or footer
            \item Add a \texttt{*} to show the header and footer bars
        \end{itemize}
    }

    % \coverframe*[crop]{image} -- same but with header/footer bars visible
    \coverframe*[p]{Earth_From_the_Perspective_of_Artemis_II.jpeg}

    \textframe{\texttt{\textbackslash coverframe*} --- with Bars}{
        The previous slide used \texttt{\textbackslash coverframe*[p]\{...\}}.
        \vspace{0.8em}
        \begin{itemize}
            \item The \texttt{*} keeps the header and footer visible over the image
            \item Useful when the section context or slide number needs to be shown
            \item The \texttt{[credit]} argument works with \texttt{*} as well
        \end{itemize}
    }

    % \coversplitframe[width]{align}{crop}{title}{image}{content}[credit]
    % width  -- textbox width (default: 0.45\paperwidth)
    % align  -- l docks box left; r docks box right
    % crop   -- l or p (same as \coverframe)
    % title  -- shown as box header with rule; pass {} to omit
    % credit -- optional; pass as final argument after content
    \coversplitframe[0.45\paperwidth]{l}{p}{Cover Split: Text Left}{Earth_From_the_Perspective_of_Artemis_II.jpeg}{
        \textbf{\texttt{\textbackslash coversplitframe[w]\{align\}\{crop\}\{title\}\{img\}\{content\}[credit]}}
        \vspace{0.5em}

        Full-bleed image with a floating translucent text box.
        \begin{itemize}
            \item \texttt{[w]} --- box width (default: \texttt{0.45\textbackslash paperwidth})
            \item \texttt{align} --- \texttt{l} docks box left; \texttt{r} docks right
            \item \texttt{title} --- box header with decorative rule; pass \texttt{\{\}} to omit
            \item \texttt{[credit]} --- optional; appended after the content argument
        \end{itemize}
    }[Credit: NASA/Artemis II]

    \coversplitframe[0.45\paperwidth]{r}{p}{}{Earth_From_the_Perspective_of_Artemis_II.jpeg}{
        Same command with \texttt{align = r} and an empty title (\texttt{\{\}}).
        \vspace{0.5em}
        \begin{itemize}
            \item Box docks to the right edge
            \item No title or rule is shown
            \item Add a \texttt{*} to show the header and footer bars
        \end{itemize}
    }

    % \coversplitframe*[width]{align}{crop}{title}{image}{content} -- with bars
    \coversplitframe*[0.45\paperwidth]{l}{p}{Cover Split with Bars}{Earth_From_the_Perspective_of_Artemis_II.jpeg}{
        \texttt{\textbackslash coversplitframe*} keeps the header and footer visible.
        \vspace{0.5em}
        \begin{itemize}
            \item Otherwise identical to \texttt{\textbackslash coversplitframe}
            \item The section name bar and footer sit on top of the image
        \end{itemize}
    }

    % ==========================================
    % DIVIDER SLIDES
    % ==========================================

    \section{Divider Slides}
    \sectionframe{Divider Slides}{\texttt{\textbackslash sectionframe} and \texttt{\textbackslash subsectionframe}}

    \textframe{\texttt{\textbackslash sectionframe\{title\}\{subtitle\}}}{
        The slide above this one was produced by \texttt{\textbackslash sectionframe}.
        \vspace{0.8em}
        \begin{itemize}
            \item \texttt{title} --- large centred section title
            \item \texttt{subtitle} --- smaller description below the rule; pass \texttt{\{\}} to omit
            \item Always a \texttt{[plain]} frame
            \item Automatically inserted at each \texttt{\textbackslash section\{\}} when the \texttt{autosec} theme option is active
        \end{itemize}
    }

    % \subsectionframe{parent}{title}
    % parent -- small context label at top (e.g. the parent section name)
    % title  -- large main subsection title
    \subsectionframe{Divider Slides}{Subsection Divider: \texttt{\textbackslash subsectionframe}}

    \textframe{\texttt{\textbackslash subsectionframe\{parent\}\{title\}}}{
        The previous slide was produced by \texttt{\textbackslash subsectionframe}.
        \vspace{0.8em}
        \begin{itemize}
            \item \texttt{parent} --- small label at top giving the parent section context
            \item \texttt{title} --- large main subsection title
            \item Useful for multi-level presentation structure
        \end{itemize}
    }

    % ==========================================
    % DARK MODE
    % ==========================================

    \section{Dark Mode}
    \sectionframe{Dark Mode}{The \texttt{darkslide} environment}

    % darkslide -- wraps any slide command(s) to render in dark mode
    \begin{darkslide}
    \textframe{The \texttt{darkslide} Environment}{
        Wrap any slide command in \texttt{\textbackslash begin\{darkslide\}...\textbackslash end\{darkslide\}}.
        \vspace{0.8em}
        \begin{itemize}
            \item Background, text, frametitle, and footer all invert to dark
            \item The \texttt{\textbackslash coversplitframe} text box colour also flips automatically
            \item Multiple commands can be nested inside one \texttt{darkslide} block
            \item The theme accent colour is preserved
            \item Compatible with \texttt{\textbackslash setbgimage}
        \end{itemize}
    }
    \end{darkslide}

    \begin{darkslide}
    \splitframe[0.5\textwidth]{r}{Dark Mode: Split Frame}{apj460815f8_hr.jpg}{
        \texttt{darkslide} works with \texttt{\textbackslash splitframe}.
        \vspace{0.5em}
        \begin{itemize}
            \item Background inverts to dark
            \item Text and bullets become white
        \end{itemize}
    }
    \end{darkslide}

    \begin{darkslide}
    \coversplitframe[0.45\paperwidth]{r}{p}{Dark Mode: Cover Split}{Earth_From_the_Perspective_of_Artemis_II.jpeg}{
        \texttt{darkslide} works with \texttt{\textbackslash coversplitframe}.
        \vspace{0.5em}
        \begin{itemize}
            \item Textbox background flips to dark gray
            \item Text flips to white
        \end{itemize}
    }
    \end{darkslide}

    % ==========================================
    % CREDITS
    % ==========================================

    \section{Credits}
    \sectionframe{Credits}{\texttt{\textbackslash credit}}

    \textframe{\texttt{\textbackslash credit\{text\}}}{
        A general-purpose credit label, placeable inside any frame body.
        \vspace{0.8em}
        \begin{itemize}
            \item Renders a small dark semi-transparent pill with white italic text at the bottom-right corner
            \item Works on both light and dark slides, and over full-bleed images
            \item Call it anywhere inside a \texttt{\textbackslash textframe} or \texttt{\textbackslash splitframe} content argument
            \item Built in to \texttt{\textbackslash imageframe}, \texttt{\textbackslash coverframe}, and \texttt{\textbackslash coversplitframe} via an optional final \texttt{[credit]} argument
        \end{itemize}
        \vspace{0.8em}
        \texttt{\textbackslash credit\{Author et al. 2025, ApJ, 123, 456\}}
        \credit{Author et al. 2025, ApJ, 123, 456}
    }

    % ==========================================
    % BACKGROUND IMAGES
    % ==========================================

    \section{Background Images}
    \sectionframe{Background Images}{\texttt{\textbackslash setbgimage} and \texttt{\textbackslash clearbgimage}}

    \textframe{\texttt{\textbackslash setbgimage} and \texttt{\textbackslash clearbgimage}}{
        Persistently set a background image for all subsequent frames.
        \vspace{0.8em}
        \begin{itemize}
            \item \texttt{\textbackslash setbgimage[crop]\{path\}\{opacity\}} --- activates the background image
            \item \texttt{[crop]} --- \texttt{l} or \texttt{p}, same as \texttt{\textbackslash coverframe} (default: \texttt{l})
            \item \texttt{opacity} --- \texttt{0.0} (invisible) to \texttt{1.0} (full); \texttt{0.1}--\texttt{0.3} is typical
            \item \texttt{\textbackslash clearbgimage} --- removes the image, restoring the global default
            \item A global default can be set via theme options: \texttt{bgimage}, \texttt{bgopacity}, \texttt{bgcrop}
        \end{itemize}
    }

    % Background image active from here
    \setbgimage[p]{Earth_From_the_Perspective_of_Artemis_II.jpeg}{0.15}
    \textframe{Background Image: Light Mode (opacity 0.15)}{
        \texttt{\textbackslash setbgimage[p]\{Earth\_...\}\{0.15\}} was called before this frame.
        \vspace{0.8em}
        \begin{itemize}
            \item The image is rendered inside \texttt{background canvas}, beneath all content
            \item Header, footer, and frame title are unaffected
            \item Low opacity keeps text legible over most images
        \end{itemize}
    }

    \setbgimage[p]{Earth_From_the_Perspective_of_Artemis_II.jpeg}{0.35}
    \textframe{Background Image: Higher Opacity (0.35)}{
        Same image, opacity raised to \texttt{0.35}.
        \vspace{0.8em}
        \begin{itemize}
            \item Image is more prominent but text readability decreases
            \item Adjust to suit the contrast of your image and content
        \end{itemize}
    }

    \setbgimage[p]{Earth_From_the_Perspective_of_Artemis_II.jpeg}{0.25}
    \begin{darkslide}
    \textframe{Background Image: Dark Mode (opacity 0.25)}{
        \texttt{\textbackslash setbgimage} and \texttt{darkslide} can be combined.
        \vspace{0.8em}
        \begin{itemize}
            \item Dark canvas renders first; image is overlaid at the specified opacity
            \item Adjust opacity to balance image visibility and text contrast
        \end{itemize}
    }
    \end{darkslide}

    % Restore clean background before closing slide
    \clearbgimage

    % ==========================================
    % CLOSING
    % ==========================================

    % \imagetitleframe with banner at the top
    \imagetitleframe[north]{p}{Earth_From_the_Perspective_of_Artemis_II.jpeg}

\end{document}
